\section{Summary}

\subsection{Results}

Table~\ref{tab:design-space-results} summarizes the explored design points.
The first three columns identify the architectural choice for each run: input
bus width, parallel input consumption, and the per-layer neuron split. The
remaining columns are the judging metrics: FMAX, throughput, latency, FPS, and
resource utilization. The blank FPS and resource fields are placeholders for
the final submission measurements once the selected point is locked.

\begin{strip}
  \centering
  \captionof{table}{Current design-space results summary.}
  \label{tab:design-space-results}
  \resizebox{\textwidth}{!}{%
    \begin{tabular}{@{}rrrrrrrr@{}}
      \toprule
      \textbf{Input Bus Width} & \textbf{Parallel Inputs} &
      \textbf{Parallel Neurons} & \textbf{FMAX} &
      \textbf{Throughput} & \textbf{Latency} &
      \textbf{FPS} & \textbf{Resource Utilization} \\
      \midrule
      64   & 32  & ---           & 722 & 105  & 290.2 & --- & --- \\
      128  & 64  & 64,32,10      & 635 & 56.1 & 195.2 & --- & --- \\
      256  & 128 & 64,64,10      & 501 & 32.2 & 130.2 & --- & --- \\
      512  & 128 & 256,64,10     & --- & 17.8 & 95.2  & --- & --- \\
      1024 & 128 & 256,64,10     & --- & 10.6 & 78.4  & --- & --- \\
      \bottomrule
    \end{tabular}%
  }
\end{strip}

\subsection{Implementation and Tradeoffs}

The implementation target was to maximize throughput first and then optimize
for clock frequency. Another design priority was parameterization. I wanted 
every component to be parameterized to be fully scalable to any power-of-two amount of 
parallel neurons, parallel inputs, and input width. This allowed me to gather 
these results for various configurations as I scaled up the input bus width.

Throughput is increased by replicating neuron processors
and the corresponding weight and threshold RAM banks, allowing more neurons to
be processed in parallel each cycle. However, that same replication creates
difficult timing paths because the RAM-side control and data signals become
large fanout nets with significant routing cost. To address this, multiple
pipeline stages were inserted on the interfaces that communicate with the
weight and threshold RAM units, breaking these long paths into shorter stages
and improving timing closure.

A similar approach is used in the compute reductions. The XNOR accumulation and
argmax logic are implemented as recursive balanced trees, which makes it
natural to pipeline the computation between tree levels. This reduces the logic
depth of each stage and allows the design to scale to wider buses and larger
parallel neuron counts without collapsing FMAX. The design is also heavily
parameterized, which made it possible to evaluate the architecture across a
range of bus widths, parallel-input counts, and neuron configurations while
preserving the same overall structure.

The design also supports functionality beyond what was tested for the project.
Notably, the design supports partial configuration and input buses as specified
by TKEEP, it also supports complete reconfiguration of weights and thresholds in
between image streams, as well as customizable hidden layers and topologies.

\subsection{Verification/Testbench}

The base testbench was extended into a UVM flow to verify functionality beyond
the nominal contest path and to provide more complete coverage of the design.
The UVM sweep was used to verify the extended functionality beyond the nominal
contest path: partial \texttt{TKEEP} configuration and input streams,
interleaved and reordered configuration traffic, weights-only and
thresholds-only reconfiguration, partial-layer updates, deeper hidden-layer
topologies, extreme threshold/weight/pixel values, long inter-transaction
gaps, output backpressure, and reset/reconfigure behavior.

\begin{itemize}
  \item Functional coverage was collected with \texttt{make coverage-sweep}.
  \item A text coverage report was generated with \texttt{make coverage-sweep-report}.
  \item The achieved-coverage report is written to\\
  \hspace*{1.5em}\texttt{coverage/bnn\_fcc\_coverage\_sweep\_test\_coverage.txt}.
  \item Supports partial \texttt{TKEEP} configuration and input streams.
  \item Supports reconfiguring weights and thresholds after a few images.
  \item Supports interleaved configuration streams, including thresholds-first and out-of-order layer updates.
  \item Supports deeper hidden-layer stacks; see \texttt{make sim-ten-hidden}.
  \item Tested long delay gaps and extreme weight, threshold, and pixel values.
  \item Scales to higher neuron counts when the input bus width increases.
\end{itemize}
